\chapter*{Tóm Tắt}
\addcontentsline{toc}{chapter}{Tóm Tắt}

Đồ án xây dựng một ứng dụng thuật toán trên mạng đường Hà Nội gồm hai chế độ có
chung nền tảng đồ thị. Chế độ thứ nhất giải bài toán lập lộ trình qua nhiều địa
điểm: Dijkstra được dùng để tính đường đi ngắn nhất giữa các điểm quan tâm,
Held--Karp tìm thứ tự ghé tối ưu chính xác cho tối đa 10 điểm, còn láng giềng gần
nhất và 2-opt đóng vai trò phương pháp so sánh. Hệ thống hỗ trợ ba kiểu kết thúc
tuyến và hai hàm mục tiêu: tổng quãng đường hoặc tổng thời gian ước tính.

Chế độ thứ hai mô phỏng trò chơi truy tìm theo lượt trên đồ thị. Người trốn biết
vị trí thật của mình, trong khi máy chỉ được quan sát tập vị trí khả dĩ, dấu vết
tiết lộ theo chu kỳ, vị trí các đội truy tìm và các lối thoát. Ba chiến thuật
\textit{đuổi dấu vết}, \textit{tìm kiếm độc lập} và \textit{phối hợp chặn lối
thoát} được xây dựng trên khoảng cách Dijkstra. Ranh giới thông tin được tách rõ
trong kiểu dữ liệu quan sát nhằm ngăn thuật toán máy đọc trực tiếp trạng thái ẩn.

Phiên bản được kiểm tra ngày 24/10/2026 vượt qua 112/112 ca kiểm thử. Graph định
tuyến đóng gói chứa 350.930 nút và 864.982 cạnh. Smoke test năm địa điểm cho kết
quả 26,276 km khi tối ưu quãng đường và 29,93 phút khi tối ưu thời gian. Benchmark
riêng cho tầng thứ tự ghé trên 24 cấu hình cho thấy Held--Karp chỉ tốn dưới 7,2 ms
ở giới hạn 10 điểm, trong khi láng giềng gần nhất lệch tới 19,53\% và 2-opt lệch
tới 3,97\% so với nghiệm tối ưu. Benchmark
trò chơi gồm 360 ván cho thấy hiệu quả chiến thuật phụ thuộc mạnh vào kiểu di
chuyển của người trốn và cấu trúc sân; do đó báo cáo không tuyên bố một heuristic
luôn vượt trội. Kết quả minh họa giá trị của việc kết hợp thuật toán chính xác,
heuristic và kiểm thử bất biến trong một ứng dụng đồ thị thực tế.

\textbf{Từ khóa:} đồ thị đường phố, Dijkstra, Held--Karp, bài toán người du lịch,
pursuit--evasion, tập niềm tin, OpenStreetMap, Streamlit.
\clearpage
